\subsection{20.01.2026}

\subsubsection{Ziel}

Ziel des Tages war es, die Klassenstruktur des Skripts
\texttt{lcdDisplay.py} zu testen und die Messwerte des
DHT11-Sensors über das Skript \texttt{main.py} auf dem
LCD-Display auszugeben.

\subsubsection{Durchgeführte Arbeiten}

Zunächst wurde die Klassenstruktur von
\texttt{lcdDisplay.py} durch das Ausführen der Main-Methode
getestet.

Anschließend wurde überprüft, ob das Skript
\texttt{DHT11.py} weiterhin korrekt funktioniert.

Danach wurde das Skript so angepasst, dass die Temperatur-
und Luftfeuchtigkeitswerte nicht mehr ausschließlich auf der
Konsole, sondern zusätzlich auf dem LCD-Display ausgegeben
werden.

Hierfür wurde die Klasse \texttt{LcdDisplay} aus
\texttt{lcdDisplay.py} in \texttt{DHT11.py} importiert und
eine Instanz der Klasse erstellt, um die Ausgabe der
Sensordaten auf dem Display zu ermöglichen.

Zusätzlich wurde das Skript \texttt{DHT11.py} um eine
Endlosschleife erweitert, damit die Sensorwerte kontinuierlich
gemessen und aktualisiert werden.

Im weiteren Verlauf wurde innerhalb von
\texttt{DHT11.py} die Klasse \texttt{dht11Sensor}
erstellt, damit die Funktionen des Sensors zukünftig aus
dem Hauptskript aufgerufen werden können.

Der bisherige Ablauf des Skripts wurde dabei als
Main-Methode innerhalb der Klasse beibehalten, um die neue
Struktur weiterhin testen zu können.

Darüber hinaus wurde das Skript \texttt{main.py} in
\texttt{TaupunktMain.py} umbenannt.
Die bisherige Main-Methode aus \texttt{DHT11.py}
wurde anschließend in \texttt{TaupunktMain.py}
übernommen.

Zum Ende des Tages wurde die Ausgabe der Sensordaten
vorübergehend wieder auf die Konsole zurückgesetzt,
um die Fehlersuche zu vereinfachen.

Zusätzlich wurde das Skript
\texttt{TaupunktBerechnung.py} erstellt.
Hierfür wurde zunächst die vorhandene
Taupunktberechnung aus dem bereitgestellten GitHub-Repository
von :contentReference[oaicite:0]{index=0}
übernommen, um diese später an das eigene Projekt anzupassen.

\subsubsection{Problem 1: Fehler im Skript \texttt{lcdDisplay.py}}

Bei der Definition der Klasse \texttt{LcdDisplay}
wurde der Doppelpunkt nach dem Klassennamen vergessen,
wodurch ein Syntaxfehler entstand.

Zusätzlich trat ein Fehler bei der Einrichtung der
I\textsuperscript{2}C-Kommunikation auf.

Das Skript wurde ursprünglich aus der JOY-PI-Anleitung
übernommen. Dabei wurde festgestellt, dass in dem in der
Anleitung abgebildeten Beispielcode die
I\textsuperscript{2}C-Adresse des LCD-Displays fehlte.

\textbf{Lösung}

Der fehlende Doppelpunkt wurde ergänzt.
Außerdem wurde die korrekte I\textsuperscript{2}C-Adresse
manuell in den Code eingetragen.

\subsubsection{Problem 2: Fehler im Skript \texttt{DHT11.py}}

Es trat das Problem auf, dass das LCD-Display nicht korrekt
initialisiert wurde, da versehentlich die Klasse selbst
anstatt einer Instanz der Klasse gespeichert wurde.

Zusätzlich wurde beim Erstellen der Klasse
\texttt{dht11} vergessen, dass bereits ein Modul mit
dem Namen \texttt{dht11} importiert wird.

Dadurch überschrieb der Klassenname den Namen des
importierten Moduls.

\textbf{Lösung}

Die Initialisierung des LCD-Displays wurde korrigiert,
indem eine Instanz der Klasse \texttt{LcdDisplay}
erstellt wurde.

Außerdem wurde die Klasse in
\texttt{dht11Sensor} umbenannt, damit sich
Klassenname und Modulname nicht mehr gegenseitig
überschreiben.

\subsubsection{Nächste Schritte}

In der nächsten Woche soll die Taupunktberechnung
für Python angepasst und in das bestehende Projekt
integriert werden.